% !TEX root = ../../cookbook.tex
\newpage
\recipe{Alonso's Shrimp Dumplings}{Makes $\sim$50 dumplings.}{}{}

\begin{multicols}{2}
	\subsection*{For the filling}
	\begin{ingredients}
		\item 500 g shrimp
		\item 3 green onions
		\item 3 garlic cloves
		\item 2 carrots
		\item 6 leaves Chinese Cabbage\footnote{Regular cabbage works too.}
		\item 2 serrano peppers
		\item 1 tbsp oyster sauce
		\item 1 tbsp fish sauce
		\item 1 tbsp soy sauce
		\item 1 big chunk of ginger
		\item 2 tbsp sesame oil
		\item Salt and pepper
		\item Love (optional)
	\end{ingredients}
	
	\subsection*{For the wrappers\footnote{Making homemade wrappers can be risky: the quantities will vary depending on the type of flour used as well as current temperature and humidity. I will exceptionally suggest buying pre-made wrappers unless you are most confident in your wrapper-making skills, because excessively dry or moist wrappers will result in the wrong integrity and elasticity, in turn making it very difficult to obtain good dumplings.}}
	\begin{ingredients}
		\item 415 g flour
		\item 110 ml water
	\end{ingredients}
	\end{multicols}

% --- Preparation ---
\textit{Start with the filling\footnote{If you are cooking with friends, have someone start working on the wrappers.}:} Chop the \textbf{shrimp} into cm-sized cubes. 
Mince the \textbf{other solid ingredients}. 
Mix in all the \textbf{wet ingredients}.

\textit{Continue with the wrappers:} mix \textbf{flour} and \textbf{water} together until well-combined, then knead to obtain a smooth dough.
Wrap the dough in plastic and let it rest for 45-60 minutes\footnote{This is important, as it will soften the dough and give it elasticity.}.
Divide the dough into 4 equal pieces: each piece will yield approximately twelve $\sim$11 g wrappers\footnote{Keep the unused dough in plastic so it won't dry.}. 

\textit{Make the dumplings:} Put $\sim$1 tablespoon of filling in the center of each wrapper, moisten the edges of the wrapper with some water to help sealing,
then fold and seal the dumpling using whichever sealing technique you prefer\footnote{There are plenty of Youtube videos showing how to properly fold dumplings.}.

Heat a frying pan on medium heat with high smoke point vegetable oil, arrange the first batch of dumplings in the pan leaving some space between them, as they will grow during cooking.
As soon as the bottoms of the dumplings have browned, add \sfrac{1}{2} cup of water mixed with 1 tsp of flour into the pan, and cover with the lid. 
Let the dumplings steam until the water has evaporated and a brown crust has formed, approximately 4 minutes\footnote{It is wiser to risk overcooking your dumplings than eating raw shrimp.}. 
Repeat until all the dumplings are done.

% --- Description ---
\begin{recipeblock}
\end{recipeblock}